\begin{resumo}[\protect\bfseries Resumo]
Este memorial apresenta a trajetória de Stéfano De Paris Carraro no curso de
Engenharia de Software da Pontifícia Universidade Católica do Rio Grande do
Sul e sua atuação nos quatro níveis da Agência Experimental de Engenharia
de Software (AGES), entre 2024/1 e 2026/2. O relato inclui os primeiros
estudos de programação, a participação em laboratórios de pesquisa e a
entrada no trabalho profissional. Na AGES I, o projeto Decola proporcionou
experiência com prototipação, componentes de interface e desenvolvimento em
equipe. Na AGES II, o Indústria Digital ampliou a atuação para banco de
dados, desenvolvimento de serviços e testes automatizados. Na AGES III,
o VanGO envolveu responsabilidades de arquitetura, revisão de código,
integrações e infraestrutura. Na AGES IV, o Hangy exigiu maior dedicação à
gestão, ao acompanhamento das tarefas e à orientação dos colegas. Para
cada projeto, são descritas as tecnologias, as atividades por \textit{sprint} e a
avaliação da própria atuação. Ao final, o memorial discute os conhecimentos
aplicados, as dificuldades de delegação e comunicação e os efeitos do uso
de inteligência artificial no aprendizado dos estudantes.

\vspace{\baselineskip}
\noindent\textbf{PALAVRAS-CHAVE:} AGES. Engenharia de Software.
Desenvolvimento de software. Gestão de projetos. Aprendizado.
\end{resumo}
